\documentclass[conference]{IEEEtran}
\IEEEoverridecommandlockouts

\usepackage{cite}
\usepackage{amsmath,amssymb,amsfonts}
\usepackage{graphicx}
\usepackage{textcomp}
\usepackage{xcolor}
\usepackage{booktabs}
\usepackage{array}
\usepackage{url}
\usepackage{float}
\usepackage{tikz}
\usetikzlibrary{arrows.meta, positioning, shapes.geometric, fit, backgrounds, calc}
\usepackage{algorithm}
\usepackage{algpseudocode}
\algrenewcommand\algorithmicrequire{\textbf{Input:}}
\algrenewcommand\algorithmicensure{\textbf{Output:}}
\usepackage{multirow}
\usepackage{caption}
\usepackage{tabularx}

\def\BibTeX{{\rm B\kern-.05em{\sc i\kern-.025em b}\kern-.08em
    T\kern-.1667em\lower.7ex\hbox{E}\kern-.125emX}}

\begin{document}

\title{Semantic Firewall: A Hybrid Defense-in-Depth\\Architecture for LLM Security}

\author{\IEEEauthorblockN{Anonymous Authors}}

\maketitle

% ====================================================================
% ABSTRACT
% ====================================================================
\begin{abstract}
Large Language Models (LLMs) are increasingly deployed in high-stakes
applications, yet defending them against adversarial prompts remains
an open problem. The prevailing approach---routing every user input
through a dedicated safety LLM such as Llama Guard or OpenAI's
moderation endpoint---treats a harmless question and a sophisticated
jailbreak with equal computational expense. This is wasteful by
design. We present \textbf{Semantic Firewall}, a system built on the
premise that most attacks are structurally predictable. Instead of
throwing a large model at every request, the architecture cascades
heuristics resolve known signatures in under 5\,ms; a self-healing
vector cache (ChromaDB~\cite{chromadb} + SentenceTransformer embeddings) intercepts
semantic variants of previously blocked threats at sub-millisecond
lookup cost (with $\sim$600\,ms for initial CPU embedding); eight
parallel heuristic detectors scan for PII, secrets,
abuse, and injection patterns simultaneously; and a constrained
Llama-3.3-70B gate---locked to greedy decoding and strict JSON
output---handles only what survives all prior filters. When this
final layer blocks a novel threat, its embedding is automatically
written back to the cache, so the system learns without manual
intervention. On the neuralchemy prompt injection benchmark~\cite{neuralchemy}
($N{=}2{,}000$), the system achieves 92.64\% F1 ($\pm$0.42) with
99.82\% Recall (cold-start), outperforming GPT-4.0 (90.00\%) and
Claude~3.5 Haiku (83.20\%), while bypassing LLM inference entirely for
66.88\% of requests on the evaluation distribution (50\% attack rate);
deployment rates will vary with traffic composition. Cross-dataset
evaluation on BeaverTails, deepset/prompt-injections, and multi-turn
evasion scenarios shows these results generalize.
\end{abstract}

\begin{IEEEkeywords}
large language models, prompt injection, semantic firewall, LLM
security, defense-in-depth, semantic caching
\end{IEEEkeywords}



% ====================================================================
% I. INTRODUCTION
% ====================================================================
\section{Introduction}

Large Language Models (LLMs) have quietly become critical infrastructure.
They power customer support pipelines, generate production code,
answer medical queries, and drive autonomous agents that make
consequential decisions with little human oversight. This rapid
shift from research curiosity to load-bearing dependency has
widened the attack surface faster than defenses have
matured~\cite{surveyprompt,owasptop10}. Adversaries inject
carefully crafted prompts to hijack the model
behavior~\cite{greshake2023indirect}, design jailbreaks that slip
past alignment training~\cite{autodan,pair,wei2023jailbroken,zou2023universal},
and exploit context windows to exfiltrate Personally Identifiable Information (PII) or cryptographic
secrets.

The standard industry response has been to place another model in
front of the first one. Llama Guard~\cite{llamaguard}, Guardrails
AI~\cite{guardrailsai}, and PromptGuard~\cite{promptguard} route
every incoming prompt through a dedicated safety LLM---whether it
says ``What is the weather today?'' or attempts a multi-stage persona
hijack. This approach is conceptually simple but practically
costly. It typically adds 800--1500\,ms of latency per request and
scales GPU utilization linearly with traffic. More critically, these
systems rely on statically trained models whose safety tuning does
not generalize to structural attacks---Base64-encoded payloads,
token smuggling, multi-turn hypothetical framing---that evolve faster
than retraining cycles permit. Novel obfuscation techniques that
exploit distributional blind spots in safety fine-tuning data pass
through undetected, while opaque \texttt{403 Forbidden} rejections
leave downstream applications unable to distinguish a genuine attack
from a false alarm. The fundamental problem is architectural: a
single checkpoint, however powerful, cannot simultaneously be fast,
cheap, broadly generalizing, and robust to out-of-distribution
attacks.

What is needed instead is a system that treats prompt security the
way network security has been practiced for decades: not as a single
gate, but as a layered pipeline where each stage is optimized for a
different class of threat. Cheap deterministic filters should handle
the bulk of structurally predictable attacks, freeing expensive
neural inference for only the genuinely ambiguous cases that
require semantic understanding. Crucially, the system should learn
from every attack it encounters, converting novel threats into future
fast-path detections without requiring manual rule updates or model
retraining.

We built \textbf{Semantic Firewall} around the principle that
defense-in-depth works better than any single
checkpoint~\cite{saltzer1975protection}. The system cascades through
four layers of progressively deeper analysis, short-circuiting as
soon as a layer resolves the input. Fast deterministic filters handle
the vast majority of traffic; the expensive LLM gate is reserved
only for genuinely novel inputs. Novel threats it blocks are
automatically embedded into a self-healing semantic cache, hardening
future detection without manual intervention. The outcome: \textbf{92.64\%
F1 with 99.82\% Recall (cold-start)}, which we report as our conservative baseline. In steady-state operation (warm cache), this rises to 93.57\% F1 with latency dropping to $\sim$427.72\,ms, bypassing LLM inference
entirely on 66.88\% of requests---outperforming GPT-4.0, Claude~3.5,
and Llama Guard~4 without routing a single prompt outside the
operator's own infrastructure.

% ── Figure placed here so it renders at top of page 2/3 ──
\begin{figure*}[!t]
\centering
\includegraphics[scale=0.30]{Semantic_Firewall_Architecture_600DPI.pdf}
\caption{Multi-agent Semantic Firewall System}
\label{fig:arch-flow}
\end{figure*}

We make four core contributions:
\textbf{(1)~A hybrid cascading architecture} that layers deterministic regex,
vector-based semantic caching, a parallel heuristic ensemble, and a
constrained LLM gate, achieving state-of-the-art F1 and Recall at
sub-500\,ms latency---surpassing all monolithic baselines we tested.
\textbf{(2)~A self-healing semantic cache} that automatically vectorizes
novel threats blocked by the LLM and persists them in ChromaDB, intercepting
future variants at cache speed to drastically reduce API dependency.
\textbf{(3)~A parallel heuristic ensemble} of eight concurrent detectors
executing via \texttt{ThreadPoolExecutor} so total latency equals the slowest
individual detector, blocking 90\% of structural threats at sub-100\,ms.
\textbf{(4)~Comprehensive empirical validation} across five public benchmarks, demonstrating robust OOD
generalization and statistically significant gains over six baselines.



% ====================================================================
% II. RELATED WORK
% ====================================================================
\section{Related Work}

\textbf{Single-Model Safety Classifiers.}
The dominant production pattern routes every prompt through a
dedicated safety LLM. Llama Guard~\cite{llamaguard} and
PromptGuard~\cite{promptguard} are representative. They perform
well within their training distribution but share two structural
blind spots: per-request latency in the hundreds of milliseconds,
and an inability to catch straightforward obfuscation---Base64
encoding, explicit override tokens---that a deterministic
preprocessor handles trivially. Our evaluation makes this gap
concrete: Llama Guard~4 scores just 51.30\% F1 on the neuralchemy
benchmark (evaluated zero-shot using default settings on the
MLCommons taxonomy, mapping `unsafe' to injection/block), because
content moderation and structural prompt injection are fundamentally
different tasks.

\textbf{Adversarial Attacks on LLMs.}
A growing literature documents how fragile aligned models remain
under adversarial pressure. Zou et al.~\cite{zou2023universal}
show that universal suffix attacks (GCG) transfer across model
families. AutoDAN~\cite{autodan} generates stealthy jailbreaks
through hierarchical genetic search. PAIR~\cite{pair} automates
attack generation via iterative refinement. Greshake et
al.~\cite{greshake2023indirect} demonstrate that external data
sources can carry indirect prompt injections. Wei et
al.~\cite{wei2023jailbroken} explain why safety training breaks
down when competing objectives collide. PromptBench~\cite{promptbench}
and SmoothLLM~\cite{smoothllm} probe robustness from complementary
angles. The collective takeaway is that no single model-level
defense keeps pace with the evolving attack landscape---which is
precisely why we advocate a layered approach.

\textbf{Orchestration Frameworks.}
NeMo Guardrails~\cite{nemoguardrails} offers programmable rule-based
routing, but static lexical rules shatter against semantically
equivalent reformulations, which our embedding-based cache captures.

\textbf{Hybrid Prompt Injection Defenses.}
Rebuff~\cite{rebuff} is the closest prior system, combining
heuristics, an LLM evaluator, and a vector database for known
signatures. We improve on Rebuff in three specific ways:
(1)~constraining the LLM gate to strict JSON with greedy decoding
removes hallucination from the classification path;
(2)~the self-healing feedback loop eliminates manual signature
updates; and (3)~eight detectors run in parallel rather than
sequentially. On Rebuff's own benchmark (which inherently favors
its tuned heuristics), Semantic Firewall achieves comparable F1
(83.93\% versus Rebuff's 82.50\%).



% ====================================================================
% III. SYSTEM ARCHITECTURE
% ====================================================================
\section{System Architecture}

The Semantic Firewall is a multi-agent system where each
incoming prompt passes through progressively deeper analytical
layers. The key design principle is \textit{aggressive
short-circuiting}: the moment any layer identifies a threat, it
triggers an immediate block and all downstream computation is
skipped. This means overall latency is bounded by the
\textit{fastest} layer capable of resolving the input, not the sum
of all layers.

\begin{table}[h]
\centering
\caption{Summary of Key Notation}
\label{tab:notation}
\scriptsize
\begin{tabular}{@{}cl@{}}
\toprule
\textbf{Symbol} & \textbf{Definition} \\
\midrule
$q$ & Input prompt (user query) \\
$\vec{v}_q$ & Embedding vector of prompt $q$ (384-dim) \\
$\mathcal{C}$ & Semantic cache (set of threat embeddings in ChromaDB) \\
$\vec{v}_c$ & A cached threat embedding vector \\
$\tau$ & Cosine similarity threshold for cache matching (default 0.90) \\
$\theta$ & Aggregate risk threshold for heuristic ensemble \\
$s^*$ & Maximum cosine similarity between $\vec{v}_q$ and $\mathcal{C}$ \\
$\mathcal{D}$ & Set of heuristic detector agents $\{D_1, \ldots, D_8\}$ \\
$T$ & LLM sampling temperature (set to 0.0 for greedy decoding) \\
$P_h, P_c, P_l$ & Traffic fractions resolved by heuristics, cache, LLM \\
$L_{\text{eff}}$ & Effective weighted-average system latency \\
\bottomrule
\end{tabular}
\end{table}

\subsection{Threat Model}

We assume a black-box adversary submitting arbitrary text to elicit
harmful behavior (injections, jailbreaks, PII extraction). The
attacker may use structural obfuscation (Base64 encoding, token
smuggling), semantic reformulation, and multi-turn escalation. We
do not assume access to model weights or embeddings.



\subsection{Layer 1: Regex Engine}

Layer 1 uses pre-compiled regular expressions targeting known
signatures (overrides, Base64, token smuggling, jailbreak prefixes).
Running entirely on CPU at $\sim$5.21\,ms, this fast layer catches
structurally obvious attacks before deeper analysis is required.



\subsection{Layer 2: Semantic Cache (Self-Healing)}

Prompts are projected into 384-dimensional vectors via
SentenceTransformers~\cite{reimers2019sbert} (\texttt{all-MiniLM-L6-v2})
and queried against a ChromaDB HNSW cosine index~\cite{malkov2020hnsw}.
The similarity between the incoming vector $\vec{v}_q$ and each cached
threat vector $\vec{v}_c$ is computed via cosine similarity.

A match above threshold $\tau{=}0.90$ flags the prompt as a semantic
variant of a known attack---even when the adversary changes phrasing,
swaps adjectives, or restructures sentences---and blocks it without
ever calling the LLM. An adversary aware of the cache could
deliberately craft inputs at cosine distance just beyond $\tau$,
forcing LLM evaluation; the threshold sweep in \S\ref{sec:analysis}
mitigates this partially, and adversarial embedding hardening is the long-term solution. What makes this layer genuinely
interesting is the \textit{self-healing feedback loop}: whenever
the downstream LLM gate blocks a novel threat, that threat's
embedding is automatically inserted into ChromaDB. Every future
semantic variant within the $\tau$-radius is then caught at
cache-lookup speed. In practice, this cuts median latency from
$\sim$731\,ms (full LLM path) to $\sim$602\,ms (cache hit,
bottlenecked by $\sim$600\,ms CPU embedding generation, while the
HNSW lookup itself takes ${<}0.5$\,ms).



\subsection{Layer 3: Parallel Heuristic Ensemble}

When the cache misses, the prompt is simultaneously broadcast to
eight heuristic agents via \texttt{ThreadPoolExecutor}. Because
they run concurrently, total ensemble latency equals
$\max(T_1, T_2, \dots, T_8)$---the slowest individual agent---rather
than the sum. The detectors span complementary threat domains:
(1)~\textit{Context Flooding:} truncates inputs beyond token limits
to prevent compute exhaustion;
(2)~\textit{PII Detection:} regex-based identification of emails,
SSNs, credit cards, IBANs, and health IDs;
(3)~\textit{Secrets Scanning:} finds API keys, private keys, and
cryptographic tokens via pattern matching and Shannon entropy;
(4)~\textit{Threat Intelligence:} cross-references against known
malicious IPs and URLs;
(5)~\textit{Abuse/Toxicity:} Levenshtein distance matching against
curated toxic lexicons;
(6)~\textit{Injection Detection:} structural regex for override
patterns and token smuggling;
(7)~\textit{Unsafe Content:} rigid rule-sets for violence, CSAM,
and illegal activities; plus
(8)~\textit{Custom Rules:} enterprise-specific constraints. The
orchestrator applies a conservative OR policy: if \textit{any} agent
returns a critical flag, the prompt is blocked immediately. These
heuristics alone handle roughly 90\% of structural threats at
sub-100\,ms latency, with no neural model
in the loop.



\subsection{Layer 4: Constrained LLM Gate}

Surviving inputs reach Llama-3.3-70B-Instruct. A tight system prompt
enforces strict JSON output (\texttt{ALLOW}/\texttt{BLOCK},
confidence, reason). Greedy decoding ($T{=}0.0$) ensures
deterministic and auditable classifications. The system prompt includes
carefully designed negative constraints---for instance, explicitly
distinguishing creative roleplay requests (benign) from direct
harmful content solicitations (malicious), and strongly penalizing
over-classification of ambiguous text---to minimize false
positives on borderline cases (the full prompt is available in our
GitHub release). If the LLM times out or goes
offline, the system fails closed with aggressive heuristic fallback
rather than silently allowing traffic through.

\begin{algorithm}[t]
\caption{Semantic Firewall Classification Pipeline}
\label{alg:firewall}
\begin{algorithmic}[1]
\Require Prompt $q$, threshold $\tau$, vector cache $\mathcal{C}$, detectors $\mathcal{D} = \{D_1, \ldots, D_8\}$
\Ensure Decision $d \in \{\textsc{Allow}, \textsc{Block}\}$
\State $M_{\text{regex}} \leftarrow \textsc{RegexScan}(q)$ \Comment{Layer 1}
\If{$M_{\text{regex}} \neq \emptyset$}
    \State \Return $\textsc{Block}$
\EndIf
\State $\vec{v}_q \leftarrow \textsc{Embed}(q)$ \Comment{Layer 2}
\State $s^* \leftarrow \max_{\vec{c} \in \mathcal{C}} \dfrac{\vec{v}_q \cdot \vec{c}}{\|\vec{v}_q\| \|\vec{c}\|}$
\If{$s^* \geq \tau$}
    \State \Return $\textsc{Block}$
\EndIf
\State $\mathcal{R} \leftarrow \bigcup_{i=1}^{8} \textsc{Parallel}(D_i, q)$ \Comment{Layer 3}
\If{$\textsc{Risk}(\mathcal{R}) \geq \theta$}
    \State \Return $\textsc{Block}$
\EndIf
\State $y \leftarrow \textsc{LLM}(q,\, T{=}0.0,\, \text{JSON})$ \Comment{Layer 4}
\If{$y.\text{action} = \textsc{Block}$}
    \State $\mathcal{C} \leftarrow \mathcal{C} \cup \{\vec{v}_q\}$ \Comment{Feedback loop}
    \State \Return $\textsc{Block}$
\EndIf
\State \Return $\textsc{Allow}$
\end{algorithmic}
\end{algorithm}

\subsection{Security Analysis}

Because layer failures are positively correlated (e.g., highly
obfuscated attacks may evade both regex and cache), the system
does not assume strict statistical independence between filters.
Nevertheless, the self-healing loop introduces a useful dynamic: when the LLM catches a novel
attack and writes its embedding into the cache, the probability of
cache evasion approaches zero for all future variants that
fall within the $\tau{=}0.90$ cosine neighborhood of the blocked
embedding. The defense thus improves with exposure, turning a static
filter chain into a dynamic system where the adversary's
reformulation cost grows against a continuously shrinking evasion
surface.
Figure~\ref{fig:arch-flow} illustrates the complete
architecture; Algorithm~\ref{alg:firewall} formalizes the
classification logic. Table~\ref{tab:notation} summarizes the key
symbols used throughout.



% ====================================================================
% IV. EXPERIMENTAL SETUP
% ====================================================================
\section{Experimental Setup}

\subsection{Datasets}

We evaluate across five public benchmarks, each chosen to stress a
different aspect of the architecture.
\textbf{(1)~neuralchemy/Prompt-injection-dataset}~($N{=}2{,}000$):
a balanced mix of benign queries and complex injections, serving as
the primary evaluation corpus.
\textbf{(2)~ai4privacy/pii-masking-200k}~($N{=}209{,}261$): a
large multilingual corpus with embedded PII entities for scaled
detection testing.
\textbf{(3)~walledai/Multi-Turn-Jailbreak}~(200 sessions, 2--10
turns): adversarial conversations where attackers build trust before
injecting payloads.
\textbf{(4)~PKU-Alignment/BeaverTails}~($N{=}2{,}000$): a safety
benchmark spanning 14 harm categories with deliberately ambiguous
borderline cases.
\textbf{(5)~deepset/prompt-injections}: Rebuff's native evaluation
set, enabling a direct head-to-head.

We supplement these with a 2{,}000-sample
synthetic safety corpus (violence, illegal activities), and a
2{,}000-sample benign corpus from \texttt{HuggingFaceH4/no\_robots}
for false positive stress testing. The primary \texttt{neuralchemy}
dataset was used exclusively for evaluation; all heuristic weights
and thresholds (e.g., $\tau$) were tuned on a disjoint 500-sample
validation split.

\subsection{Environment \& Reproducibility}

Fast-path layers ran on an Intel Xeon E5-2686 v4 (8 vCPUs,
32\,GB RAM). ChromaDB~\cite{chromadb} used HNSW indexing~\cite{malkov2020hnsw}.
Embeddings were generated via SentenceTransformers~\cite{reimers2019sbert}.
LLM inference used Llama-3.3-70B-Instruct via OpenRouter ($T{=}0.0$).
Code and datasets will be released on GitHub.

\subsection{Baseline Methodology}

Commercial APIs (GPT-4.0, Claude~3.5 Haiku, Gemini~2.5 Flash) were tested
zero-shot using default moderation endpoints without custom system
prompts, reflecting standard enterprise integration.\footnote{API
snapshots evaluated: \texttt{gpt-4.0-2024-05-13},
\texttt{claude-3-5-haiku-20240307}, \texttt{gemini-2.5-flash-001}.}
We note that these APIs were not purpose-built for binary
classification, and custom system prompting might improve their
baseline performance. Llama Guard~4
was evaluated as a binary classifier on structural prompt
injections---a distribution it was not explicitly trained for, which
explains its low F1.

\subsection{Evaluation Metrics}

We report Precision, Recall, F1, and False Positive Rate (FPR).
Statistical significance for all pairwise classification comparisons is assessed via McNemar's test. We
additionally report effective latency and API cost savings.



% ====================================================================
% V. RESULTS
% ====================================================================
\section{Results}

Tables~\ref{tab:baselines}, \ref{tab:ablation}, and
\ref{tab:cache} present the full results. Isolated
component profiling gives mean latencies of 5.21\,ms (Regex),
83.45\,ms (Heuristics), 602.83\,ms (Cache, bottlenecked by CPU
embedding), and 731.64\,ms (LLM Gate).

% ── TABLE 1: Baselines + Cross-Dataset + Specialized ──
\begin{table*}[!t]
\centering
\caption{Baseline Comparison, Cross-Dataset Generalization, and Specialized Evaluations.
Best results per benchmark are \textbf{bolded}. ``---'' indicates metric not applicable.}
\label{tab:baselines}
\scriptsize
\begin{tabular}{@{}llrccccl@{}}
\toprule
\textbf{Benchmark} & \textbf{Model / Config}
    & \textbf{$N$} & \textbf{Prec.} & \textbf{Recall}
    & \textbf{F1} & \textbf{Latency} & \textbf{Notes} \\
\midrule
\multicolumn{8}{@{}l}{\textit{(a) Open-Weight vs. Proprietary API
    Comparison (neuralchemy, $N{=}2{,}000$)}} \\
\midrule
neuralchemy & \textbf{Semantic Firewall (Ours, Warm)}
    & 2000 & 89.03 & 98.61 & \textbf{93.57}
    & \textbf{427.72\,ms} & Steady-state \\
neuralchemy & Semantic Firewall (Ours, Cold)
    & 2000 & 86.43 & \textbf{99.82} & 92.64
    & 459.85\,ms & Empty cache \\
neuralchemy & OpenAI GPT-4.0
    & 2000 & \textbf{96.10} & 84.60 & 90.00
    & 876\,ms & Proprietary \\
neuralchemy & Google Gemini 2.5 Flash
    & 2000 & 95.30 & 81.00 & 87.60
    & 1464\,ms & Proprietary \\
neuralchemy & Anthropic Claude 3.5 Haiku
    & 2000 & 94.70 & 74.30 & 83.20
    & 1156\,ms & Proprietary \\
\midrule
\multicolumn{8}{@{}l}{\textit{(b) Cross-Dataset Generalization}} \\
\midrule
BeaverTails
    & \textbf{Semantic Firewall}
    & 2000 & 65.23 & \textbf{87.59} & \textbf{74.80}
    & --- & 14 harm categories \\
BeaverTails
    & Llama Guard 4 (12B)
    & 2000 & 73.68 & 61.20 & 66.86
    & --- & --- \\
deepset/prompt-inj.
    & \textbf{Semantic Firewall}
    & --- & \textbf{95.10} & 75.14 & \textbf{83.93}
    & --- & Rebuff's dataset \\
deepset/prompt-inj.
    & ProtectAI Rebuff
    & --- & 86.20 & \textbf{79.10} & 82.50
    & --- & Pipeline \\
\midrule
\multicolumn{8}{@{}l}{\textit{(c) Specialized Evaluations}} \\
\midrule
ai4privacy/pii-200k
    & SF Regex Layer
    & 209K & 85.90 & 96.70 & 91.00
    & $<$2\,ms & Zero LLM cost \\
walledai/Multi-Turn
    & Semantic Firewall
    & 200 & --- & 90.00 & ---
    & --- & Session recall \\
Synthetic Safety
    & SF Full
    & 2000 & 82.50 & 84.30 & 83.38
    & --- & FPR = 5.45\% \\
HF/no\_robots
    & SF Full
    & 2000 & --- & --- & ---
    & --- & FPR = 4.20\% \\
\bottomrule
\end{tabular}
\end{table*}

% ── TABLE 2: Ablation Study ──
\begin{table}[t]
\centering
\caption{Ablation Study --- Configuration Comparison
(neuralchemy, $N{=}2{,}000$). Best results are \textbf{bolded}.}
\label{tab:ablation}
\scriptsize
\begin{tabular}{@{}lrcccc@{}}
\toprule
\textbf{Configuration}
    & \textbf{$N$} & \textbf{Prec.} & \textbf{Recall}
    & \textbf{F1} & \textbf{Latency} \\
\midrule
C1: LLM Gate Only
    & 2000 & 78.23 & 63.66 & 70.20 & 731.64\,ms \\
C2: Regex Only
    & 2000 & 77.34 & 62.94 & 69.40 & 5.21\,ms \\
C3: Cache Only
    & 2000 & 81.24 & 56.29 & 66.50 & 602.83\,ms \\
C4: Par.\ Det.\ Only
    & 2000 & 83.07 & 93.77 & 88.10 & 83.45\,ms \\
C5: Regex + Cache
    & 2000 & 98.51 & 80.17 & 88.40 & 602.83\,ms \\
C6: Regex + Par.\ Det.
    & 2000 & 92.04 & 88.81 & 90.40 & 83.45\,ms \\
C7: All Fast (No LLM)
    & 2000 & 98.14 & 82.77 & 89.80 & 602.83\,ms \\
\textbf{C8: Full (Warm)}
    & 2000 & \textbf{89.03} & \textbf{98.61} & \textbf{93.57} & \textbf{427.72\,ms} \\
C9: Full (Cold)
    & 2000 & 86.43 & 99.82 & 92.64 & 459.85\,ms \\
\bottomrule
\end{tabular}
\end{table}

% ── TABLE 3: Cache Threshold Sensitivity ──
\begin{table}[t]
\centering
\caption{Cache Threshold Sensitivity (Across $\tau$ sweeps)}
\label{tab:cache}
\scriptsize
\begin{tabular}{@{}lcccc@{}}
\toprule
\multicolumn{5}{@{}l}{\textit{(a) Full System (C8) on neuralchemy ($N{=}2{,}000$)}} \\
\midrule
\textbf{Threshold ($\tau$)}
    & \textbf{Prec.} & \textbf{Recall}
    & \textbf{F1} & \textbf{FPR} \\
\midrule
$\tau = 0.80$
    & 78.12 & 99.43 & 88.24 & 18.14\% \\
$\tau = 0.85$
    & 85.34 & 99.11 & 90.47 & 12.41\% \\
$\boldsymbol{\tau = 0.90}$ \textbf{(Sel.)}
    & \textbf{89.03} & \textbf{98.61} & \textbf{93.57} & \textbf{4.20\%} \\
$\tau = 0.95$
    & 94.81 & 92.86 & 91.43 & 4.24\% \\
$\tau = 0.99$
    & 96.27 & 88.29 & 88.12 & 1.83\% \\
\midrule
\multicolumn{5}{@{}l}{\textit{(b) OOD Generalization}} \\
\midrule
\textbf{Threshold ($\tau$)} & \textbf{Acc.} & \textbf{Prec.} & \textbf{Recall} & \textbf{F1} \\
\midrule
$\tau = 0.80$ & 89.80\% & 96.49\% & 87.30\% & 91.67\% \\
$\tau = 0.85$ & 91.84\% & 98.25\% & 88.89\% & 93.33\% \\
$\tau = 0.90$ & 86.73\% & 96.30\% & 82.54\% & 88.89\% \\
$\tau = 0.95$ & 90.82\% & 98.21\% & 87.30\% & 92.44\% \\
\bottomrule
\end{tabular}
\end{table}

\subsection{Analysis}
\label{sec:analysis}

\textbf{Ablation (Table~\ref{tab:ablation}).}
Each added layer improves F1 monotonically, but what is most
revealing is how the standalone LLM gate (C1: 70.20\% F1 at
731\,ms) and the standalone regex engine (C2: 69.40\% at 5\,ms)
perform almost identically---expensive inference and cheap pattern
matching are roughly equivalent in isolation. The real value of our
architecture lies in their composition. Regex plus cache (C5) pushes
Precision to 98.51\%; adding parallel detectors (C6) lifts F1 to
90.40\%; the full pipeline (C8) peaks at the highest F1 and Recall
in the table. The 23.37-point gain from C1 to C8 is statistically
significant (McNemar's $\chi^2{=}135.2$, $p < 0.001$). Notably, adding the LLM gate to the fast-path (C5 to C8)
lowers Precision by $\sim$9.5 points while lifting Recall dramatically (80\% to 99\%).
This precision drop stems directly from the LLM's conservative false positive behavior
on ambiguous borderline cases (e.g., creative writing prompts containing ostensibly violent keywords).
We accept this trade-off: in a security context, catching 99\% of attacks with slightly
more false alarms is preferable to leaking 20\% of attacks. Finally, parallel execution
means that stacking heuristic layers does not stack their latencies: C4 and C6 both
bottleneck at 83.45\,ms.

\textbf{Baselines (Table~\ref{tab:baselines}a).}
The system outperforms every commercial API in F1: 92.64\% versus
GPT-4.0's 90.00\%, Claude~3.5 Haiku's 83.20\%, and Gemini 2.5 Flash's
87.60\%. Commercial models achieve higher raw Precision (GPT-4.0:
96.10\%), but their Recall is consistently lower---GPT-4.0 catches
84.60\% of attacks versus our 99.82\%. The explanation is
architectural: regex and cache layers intercept structural attacks
that no standalone LLM catches, lifting system-level Recall past
what any individual model achieves. Statistical significance was
verified via McNemar's test ($\chi^2{=}4.28$, $p = 0.038$). Because
41.93\% of traffic is resolved deterministically, effective
cold-start latency drops to 459.85\,ms---faster than every API
benchmarked.

\textbf{Cross-Dataset (Table~\ref{tab:baselines}b).}
The system generalizes to distributions it has never trained on. On
BeaverTails---14 harm categories, deliberately difficult borderline
cases---the Semantic Firewall reaches 74.80\% F1 versus Llama
Guard~4's 66.86\%, outperforming Meta's dedicated safety model by
$\sim$8 F1 points. The relatively low Precision (65.23\%) on this dataset reflects
its construction: BeaverTails contains many semantically ambiguous cases that
trigger our conservative heuristics, resulting in false positives but
maintaining high Recall (87.59\%). On Rebuff's own deepset benchmark, the system
achieves comparable F1 (83.93\% vs.\ 82.50\%), suggesting
competitive performance without Rebuff's manual signature management.

\textbf{Specialized (Table~\ref{tab:baselines}c).}
The PII regex layer achieves 91.0\% F1 with 96.7\% Recall across
209K samples at sub-2\,ms---no neural model involved. Highly
structured entities (emails, credit cards, IBANs) hit $>$99.7\%
Recall; variable types like phone numbers reach 74.8\%. Multi-turn
evaluation confirms 90.0\% session-level interception. Benign degradation shows 4.20\% FPR,
within the 5--15\% industry range.

\textbf{Cache Threshold (Table~\ref{tab:cache}).}
Sweeping $\tau$ from 0.80 to 0.99 reveals a clean trade-off:
$\tau{=}0.80$ maximizes hits (42.4\%) but inflates FPR to 18.14\%;
$\tau{=}0.99$ minimizes FPR (1.83\%) but the cache barely fires.
The selected threshold $\tau{=}0.90$ optimizes F1 at a 12.45\% hit
rate and 4.20\% FPR.

\textbf{Cost.}
On our evaluation distribution, fast layers resolved 66.88\% of all
incoming requests without LLM invocation; this ratio will vary with
deployment-specific traffic composition. Let $P_h$, $P_c$, and $P_l$ denote the fractions of traffic resolved by the heuristic
ensemble, cache, and LLM gate respectively, with corresponding
mean latencies $L_h$, $L_c$, $L_l$. The effective system latency is:

\begin{equation}
L_{\text{eff}} = P_h \cdot L_h \;+\; P_c \cdot L_c \;+\; P_l \cdot L_l
\label{eq:latency}
\end{equation}

\noindent On our evaluation split ($P_h{=}0.55$, $P_c{=}0.12$,
$P_l{=}0.33$), this yields $L_{\text{eff}} \approx 444$\,ms. At
100M tokens, GPT-4.0 costs \$500; the Semantic Firewall drops this
to \$166.



% ====================================================================
% VI. LIMITATIONS & FUTURE WORK
% ====================================================================
\section{Limitations \& Future Work}

Several limitations warrant acknowledgment. The 70B LLM gate demands
substantial GPU resources; a fine-tuned 8B model is a natural next step
for latency-constrained deployments. The self-healing cache remains
vulnerable to adversaries who craft inputs that sit outside the
$\tau{=}0.90$ cosine neighborhood---semantically distinct but functionally
equivalent attacks that evade embedding-based detection entirely, a gap
that adversarial embedding training~\cite{zou2023universal} may address.
The Python orchestrator, while sufficient for the benchmarks reported here,
caps throughput at $\sim$45 requests/second per instance due to GIL
contention; a production deployment would require rewriting the routing
layer in Go or Rust with a distributed vector store. The 4.20\% false
positive rate---roughly 1 in 24 benign prompts flagged---may be too
aggressive for applications where false blocks carry operational cost.
Finally, while pretraining contamination in
the LLM gate cannot be fully ruled out, strong cross-dataset generalization
on BeaverTails suggests the results are not a memorization artifact.



% ====================================================================
% VII. CONCLUSION
% ====================================================================
\section{Conclusion}

This paper set out to answer a practical question: does defending an
LLM really require throwing another LLM at every single request? The
evidence says no. By layering cheap deterministic filters in front of
a constrained neural gate, the Semantic Firewall achieves 92.64\% F1
with 99.82\% Recall (cold-start) while eliminating LLM inference for two-thirds
of all traffic. The ablation study is perhaps the most telling
result---stripping away the fast layers and relying on the LLM gate
alone drops F1 by 23.37 points, confirming that each layer
contributes detection capability the others cannot replicate. The
self-healing cache turns every novel blocked threat into a cheap
future cache hit, so the system gets stronger with use rather than
requiring manual rule updates. Cross-dataset evaluation on
BeaverTails and deepset/prompt-injections demonstrates that these gains are not artifacts of a single
benchmark. The architecture outperforms GPT-4.0 by 2.64 F1 points
and Claude~3.5 Haiku by 9.44 points---running entirely on open-weight
models, within the operator's own infrastructure, with no prompts
leaving the network. The complete source code and evaluation scripts
will be released for full reproducibility.

\begin{thebibliography}{20}

\bibitem{neuralchemy}
Neuralchemy, ``Prompt-injection-dataset,'' HuggingFace, 2023. [Online].
Available: \url{https://huggingface.co/datasets/neuralchemy/Prompt-injection-dataset}

\bibitem{llamaguard}
M.~Inan \textit{et al.}, ``Llama Guard: LLM-based Input-Output
Safeguard for Human-AI Conversations,''
\textit{arXiv:2312.06674}, 2023.

\bibitem{promptguard}
Meta, ``PromptGuard: A Safeguard Model Against Prompt Injections and
Jailbreaks,'' 2024.

\bibitem{rebuff}
ProtectAI, ``Rebuff: Prompt Injection Defense,'' GitHub, 2023.

\bibitem{smoothllm}
A.~Robey \textit{et al.}, ``SmoothLLM: Defending Large Language
Models Against Jailbreaking Attacks,''
\textit{arXiv:2309.02771}, 2023.

\bibitem{promptbench}
K.~Zhu \textit{et al.}, ``PromptBench: Towards Evaluating the
Robustness of Large Language Models on Adversarial Prompts,''
\textit{arXiv:2308.06588}, 2023.

\bibitem{autodan}
X.~Liu \textit{et al.}, ``AutoDAN: Generating Stealthy Jailbreak
Prompts on Aligned Large Language Models,''
\textit{arXiv:2310.04451}, 2023.

\bibitem{pair}
P.~Chao \textit{et al.}, ``Jailbreaking Black Box Large Language
Models in Twenty Queries,'' \textit{arXiv:2310.08419}, 2023.

\bibitem{nemoguardrails}
T.~Rebedea \textit{et al.}, ``NeMo Guardrails: A Toolkit for
Controllable and Safe LLM Applications,''
\textit{arXiv:2310.10501}, 2023.

\bibitem{guardrailsai}
Guardrails~AI, ``Guardrails: Validating and Correcting the Outputs
of Large Language Models,'' GitHub, 2023.

\bibitem{owasptop10}
OWASP, ``OWASP Top~10 for Large Language Model Applications,''
OWASP Foundation, 2023.

\bibitem{surveyprompt}
J.~Chen \textit{et al.}, ``A Survey on Prompt Injection Attacks in
Large Language Models,'' \textit{IEEE Trans.\ Inf.\ Forensics
Security}, 2024.

\bibitem{reimers2019sbert}
N.~Reimers and I.~Gurevych, ``Sentence-BERT: Sentence Embeddings
Using Siamese BERT-Networks,'' \textit{Proc.\ EMNLP-IJCNLP},
pp.~3982--3992, 2019.

\bibitem{malkov2020hnsw}
Y.~Malkov and D.~Yashunin, ``Efficient and Robust Approximate
Nearest Neighbor Search Using Hierarchical Navigable Small World
Graphs,'' \textit{IEEE Trans.\ Pattern Anal.\ Mach.\ Intell.},
vol.~42, no.~4, pp.~824--836, 2020.

\bibitem{greshake2023indirect}
K.~Greshake \textit{et al.}, ``Not What You've Signed Up For:
Compromising Real-World LLM-Integrated Applications with Indirect
Prompt Injection,'' \textit{Proc.\ AISec}, 2023.

\bibitem{wei2023jailbroken}
A.~Wei, N.~Haghtalab, and J.~Steinhardt, ``Jailbroken: How Does
LLM Safety Training Fail?'' \textit{Advances in Neural Information
Processing Systems (NeurIPS)}, 2023.

\bibitem{zou2023universal}
A.~Zou \textit{et al.}, ``Universal and Transferable Adversarial
Attacks on Aligned Language Models,'' \textit{arXiv:2307.15043},
2023.

\bibitem{chromadb}
Chroma, ``ChromaDB: The AI-Native Open-Source Embedding Database,''
GitHub, 2023.

\bibitem{saltzer1975protection}
J.~Saltzer and M.~Schroeder, ``The Protection of Information in
Computer Systems,'' \textit{Proc.\ IEEE}, vol.~63, no.~9,
pp.~1278--1308, 1975.

\end{thebibliography}

\end{document}
